% !TEX root = ../report.tex

\chapter{Теоретичні основи колоризації зображень}
    \label{chap:review_theory}
        У цьому розділі розглядаються основні теоретичні аспекти колоризації зображень, включно з історією розвитку методів, класичними алгоритмічними підходами та сучасними нейромережевими моделями.
        Особлива увага приділяється різноманітним підходам: від архітектур на основі згорткової нейронної мережі (англ. Convolutional Neural Network, CNN) до дифузійної моделі (англ. Diffusion Model), а також методам оцінки якості отриманих результатів.

        Основна мета цього розділу -- проаналізувати еволюцію методів колоризації від класичних CNN до сучасних архітектур типу трансформер (англ. Transformer),
        виявити їхні обмеження щодо обчислювальної складності та збереження просторових контурів об'єктів,
        і на основі цього обґрунтувати необхідність переходу до гібридних ієрархічних архітектур на базі селективної моделі простору станів (англ. Selective State Space Model, SSM).

    \section{Поняття та історія колоризації}

        Колоризація -- це процес додавання кольорової інформації до монохромних зображень.
        З математичної точки зору, цю задачу можна сформулювати як відображення одноканального сірого зображення у триканальне кольорове зображення, наприклад, у просторі RGB, CIELAB або YUV \cite{2025-survey}.
        Оскільки одному значенню яскравості може відповідати безліч кольорових відтінків, задача є невизначеною та має характер <<один до багатьох>>, що вимагає використання апріорних знань про семантику сцени.

        Історично методи колоризації можна поділити на три основні етапи:
        \begin{enumerate}
            \item \textbf{ручна та напівавтоматична колоризація:} вимагала значних зусиль художників або використання простих евристик на основі сегментації;
            \item \textbf{методи на основі глибокого навчання:} поява великих наборів даних, таких як ImageNet та COCO, дозволила тренувати CNN для передбачення кольору \cite{2016-deep-learning}.
                    Знаковим став підхід \textit{Colorful Image Colorization} \cite{2016-colorful-image-colorization}, де задачу було переформульовано з класичної регресії на класифікацію.
                    Замість того, щоб мінімізувати усереднену похибку, яка призводить до тьмяних кольорів, модель передбачає розподіл ймовірностей для кожного пікселя по палітрі відтінків.
                    Подальший розвиток цього етапу ознаменувався впровадженням генеративно-змагальної мережі (англ. Generative Adversarial Network, GAN), що дозволило подолати проблему розмитості та суттєво підвищити фотореалістичність результатів завдяки змагальному процесу навчання;
            \item \textbf{генеративні та послідовні моделі:} сучасний етап,
            що характеризується використанням механізмів глобальної уваги трансформерів та дифузійних моделей для генерації високоякісних текстур,
            а також впровадженням SSM для лінійного моделювання контексту зображень та забезпечення високої семантичної узгодженості \cite{2023-diffusion-models-survey, 2023-mamba}.
        \end{enumerate}

        Розвиток методів глибокого навчання, зокрема поява архітектур U-Net та згаданих раніше GAN, дозволив перейти від простого розфарбовування до генерації фотореалістичних зображень \cite{2017-image-to-image-translation}. 
        Останні дослідження все частіше фокусуються на дифузійних моделях, які демонструють передові результати у задачах синтезу зображень \cite{2022-palette-image-to-image-diffusion-models}. 
        
        Поряд із цим, активного розвитку набули архітектури на базі механізмів глобальної уваги.
        Проте висока обчислювальна складність класичних трансформерів, а також обмеження лінійних та інших ефективних апроксимацій у точному моделюванні глобального контексту \cite{2021-long-range-arena, 2022-efficient-transformers-survey},
        роблять їх менш придатними для інтерактивних задач.

        Водночас новітні архітектури на базі SSM здатні забезпечити глобальне рецептивне поле з лінійною складністю.
        Хоча вони й почали з'являтися у вузькоспеціалізованих задачах відновлення інфрачервоних \cite{2024-colormamba, 2026-haqagen} та супутникових знімків \cite{2026-mfmamba}, вони залишаються недослідженими в контексті інтерактивної колоризації фотографій.
        Відсутність рішень, які б поєднували лінійну обчислювальну ефективність SSM із точним просторовим контролем за допомогою користувацьких підказок, створює простір для наукового пошуку та підкреслює актуальність запропонованого в цій роботі гібридного підходу.
    
    \section{Підходи до колоризації}

        Підходи до колоризації базуються на різних архітектурах нейронних мереж.
        Опустимо гібридні архітектури та розглянемо деякі з базових.

        \subsection{Згорткові нейронні мережі}
            Ранні підходи використовували автоенкодери для прямого передбачення каналів кольоровості (наприклад, $a$ та $b$ у просторі CIE Lab).

            Ларсссон та ін. \cite{2016-learning-representations-for-automatic-colorization, 2016-colorful-image-colorization} запропонували використовувати ознаки, вивчені для класифікації об'єктів, для задачі колоризації, демонструючи тісний зв'язок між семантикою об'єкта та його кольором.

            Інший популярний метод, \textit{Deep Koalarization} \cite{2017-deep-koalarization}, використовує Inception-ResNet-v2 для вилучення ознак та подвійний декодер для відновлення кольору.

        \subsection{Генеративно-змагальні мережі}
            З метою усунення ефекту недостатньої насиченості кольорів були запропоновані різні GAN.

            \begin{itemize}[label=$\bullet$] 
                \item Модель \textit{Pix2Pix} \cite{2017-image-to-image-translation} використовує умовний GAN (англ. Conditional Generative Adversarial Network, cGAN) для перетворення зображень, що дозволяє генерувати більш чіткі та реалістичні результати.
                \item Метод \textit{ChromaGAN} \cite{2020-chromagan} поєднує геометричну інформацію з семантичним розподілом класів для покращення правдоподібності розфарбовування.
                \item Також існують підходи для непарного навчання, такі як CycleGAN \cite{2017-unpaired-image-to-image-translation}, що дозволяють тренувати моделі без наявності пар <<чорно-біле -- кольорове>>.
            \end{itemize}

        \subsection{Трансформери}
            З появою механізму уваги дослідники почали застосовувати трансформери для колоризації.

            Наприклад, \textit{Colorization Transformer} \cite{2021-colorization-transformer} та \textit{UniColor} \cite{2022-unicolor} використовують архітектуру трансформерів для моделювання довгих залежностей у зображенні, що покращує узгодженість кольорів на великих відстанях.

            Як ще один приклад: \textit{DDColor} \cite{2023-ddcolor} використовує подвійні декодери та запити кольору (англ. Color Queries) для досягнення фотореалістичності.

        \subsection{Дифузійні моделі}
            Найбільш перспективним напрямком на сьогодні є дифузійні моделі.
            Принцип їхньої роботи полягає в ітеративному оберненні процесу деградації зображення для відновлення його структури.
            Хоча у класичних підходах під деградацією розуміють поступове додавання випадкового гаусового шуму, сучасні методи розширюють це поняття на будь-які типи спотворень.
            \begin{itemize}[label=$\bullet$]
                \item \textit{Palette} \cite{2022-palette-image-to-image-diffusion-models} -- універсальний фреймворк для перетворення зображень, який перевершує GAN та регресійні моделі без необхідності налаштування під конкретну задачу.
                \item \textit{Cold Diffusion} \cite{2023-cold-diffusion} пропонує підхід до інверсії довільних перетворень зображень без використання гаусового шуму, що відкриває нові можливості для відновлення зображень.
            \end{itemize}

            % Для відео ключовою проблемою є часова узгодженість (temporal consistency).
            % Сучасні методи, такі як \textit{Stable Video Diffusion} \cite{2023-stable-video-diffusion} та \textit{VanGogh} \cite{2025-vangogh}, адаптують дифузійні моделі для обробки послідовностей кадрів (додають вимір часу, щоб бачити відео цілком).
            % Інші підходи, наприклад \cite{2019-deep-exemplar-based-video-colorization} та \cite{2024-temporally-consistent-video-colorization}, використовують поширення ознак (feature propagation) (механічно переносить кольори з попередніх кадрів, щоб вони не змінювалися) та саморегуляризацію (метод навчання, який штрафує мережу, якщо вона змінює колір одного й того ж об'єкта у сусідніх кадрах) для уникнення мерехтіння.

        \subsection{Селективні моделі простору станів}
            Новітнім трендом у глибокому навчанні є впровадження SSM, зокрема архітектури Mamba \cite{2023-mamba}.
            Вони розроблені як обчислювально ефективна альтернатива архітектурам типу трансформер. 
            Головною перевагою SSM є здатність моделювати наскрізні глобальні залежності із суворо лінійною складністю $\mathcal{O}(N)$ відносно довжини послідовності, що є критичним фактором для забезпечення інтерактивності та роботи систем у реальному часі.

    \section{Інтерактивна колоризація}  %та багатоваріантні результати

        Оскільки колоризація є відносно суб'єктивною задачею, важливу роль відіграють методи, що дозволяють користувачеві впливати на результат:
        \begin{enumerate}
            \item \textbf{колоризація на основі прикладів:}
                методи, такі як \textit{Deep Exemplar-based Colorization} \cite{2018-deep-exemplar-based-colorization}, використовують референсне кольорове зображення для перенесення колірної палітри на цільове чорно-біле;

            \item \textbf{колоризація на основі тексту:}
                з розвитком мультимодальних моделей, таких як CLIP, з'явилася можливість керувати кольором за допомогою текстових описів,
                зокрема \textit{L-CAD} \cite{2023-l-cad} використовує дифузійні пріори для колоризації на основі описів будь-якого рівня деталізації,
                а \textit{Diffusing Colors} \cite{2023-diffusing-colors} пропонує фреймворк, де текстові підказки керують процесом дифузії, забезпечуючи баланс між автоматизацією та контролем;

            \item \textbf{колоризація на основі підказок користувача:}
                методи \textit{Real-Time User-Guided Image Colorization} \cite{2017-real-time-user-guided-image-colorization} та \textit{iColoriT} \cite{2023-icolorit} дозволяють користувачеві ставити кольорові точки або штрихи, які мережа поширює на відповідні семантичні регіони.
        \end{enumerate}
        
        Сучасна тенденція, як показано в \textit{Control Color} \cite{2025-control-color} та \textit{UniColor} \cite{2022-unicolor}, полягає у створенні уніфікованих фреймворків, що підтримують різні типи умов, такі як текст, зразок, точки, штрихи, в одній моделі.

    \section{Метрики якості колоризації}

        Оцінка якості колоризації є нетривіальним завданням через фундаментальну неоднозначність розв'язку:
        одному чорно-білому вхідному сигналу відповідає ціла множина семантично та візуально коректних колірних варіацій.
        Оскільки відображення має характер $1 \to N$, використання лише одного критерію не дозволяє повноцінно оцінити роботу моделі.
        Для комплексного аналізу розроблених архітектур у роботі застосовується комбінований набір метрик, що охоплює такі аспекти:

        \begin{itemize}[label=$\bullet$]
            \item \textbf{метрики локальної та глобальної похибки:}
                у задачах інтерактивної колоризації похибку доцільно розділяти за областю обчислення.
                Зокрема, метрика середньоквадратичної похибки (англ. Mean Squared Error, $\mathrm{MSE}$) 
                та середньоабсолютної похибки (англ. Mean Absolute Error, $\mathrm{MAE}$)
                оцінюють загальне відхилення кольору на всьому зображенні.
                Вони обчислюються за формулами:
                \begin{equation*}
                    \mathrm{MSE}(x, y) = \frac{1}{N} \sum_{i=1}^N (x_i - y_i)^2, \quad \mathrm{MAE}(x, y) = \frac{1}{N} \sum_{i=1}^N |x_i - y_i|,
                \end{equation*}
                де $x_i$ та $y_i$ -- значення $i$-го пікселя згенерованого та еталонного зображень відповідно; $N$ -- загальна кількість пікселів.
                
                Спеціалізовані метрики $\mathrm{MSE}_{\mathrm{HINT}}$ та $\mathrm{MAE}_{\mathrm{HINT}}$ розраховуються за тими ж математичними формулами, але виключно в просторових координатах наданих користувачем підказок.
                Наближення $\mathrm{MSE}_{\mathrm{HINT}}$ та $\mathrm{MAE}_{\mathrm{HINT}}$ до нуля є індикатором того,
                що мережа не ігнорує користувацький ввід, а ідеально інтегрує його у фінальний результат;

            \item \textbf{піксельні та структурні метрики:}
                пікове відношення сигналу до шуму (англ. Peak Signal-to-Noise Ratio, $\mathrm{PSNR}$) та індекс структурної схожості (англ. Structural Similarity Index, $\mathrm{SSIM}$) -- класичні методи порівняння,
                які вимірюють пряме відхилення згенерованого зображення від еталонного кольорового зображення.
                $\mathrm{PSNR}$ обчислюється як:
                \begin{equation*}
                    \mathrm{PSNR}(I, J) = 10 \log_{10} \left( \frac{\mathrm{MAX}^2}{\mathrm{MSE}(I, J)} \right),
                \end{equation*}
                де $\mathrm{MAX}$ -- максимальне можливе значення пікселя зображення.
                Метрика $\mathrm{SSIM}$ розраховується за формулою:
                \begin{equation*}
                    \mathrm{SSIM}(x, y) = \frac{(2\mu_x\mu_y + c_1)(2\sigma_{xy} + c_2)}{(\mu_x^2 + \mu_y^2 + c_1)(\sigma_x^2 + \sigma_y^2 + c_2)},
                \end{equation*}
                де $\mu_x, \mu_y$ -- локальні середні значення яскравості зображень; $\sigma_x^2, \sigma_y^2$ -- дисперсії; $\sigma_{xy}$ -- коваріація; $c_1, c_2$ -- малі константи для забезпечення чисельної стабільності.
                
                Хоча ці метрики зручні для математичного аналізу, у задачах колоризації їхня репрезентативність є обмеженою:
                якщо модель згенерує семантично правильну червону машину замість оригінальної синьої,
                вони зафіксують значну помилку, незважаючи на високу візуальну правдоподібність результату;

            \item \textbf{перцептивні генеративні метрики:}
                перцептивна метрика на базі глибоких ознак (англ. Learned Perceptual Image Patch Similarity, $\mathrm{LPIPS}$) використовується для оцінки попарної схожості зображень на рівні глибоких ознак нейронної мережі VGG:
                \begin{equation*}
                    \mathrm{LPIPS}(x, x_0) = \sum_l \frac{1}{H_l W_l} \sum_{h,w} \| w_l \odot (\hat{y}^l_{hw} - \hat{y}^l_{0hw}) \|_2^2,
                \end{equation*}
                де $\hat{y}^l, \hat{y}^l_0$ -- карти ознак $l$-го шару для згенерованого та еталонного зображень; $H_l, W_l$ -- просторові розміри цих карт; $w_l$ -- вектор ваг; $\odot$ -- поелементне множення.
                
                Натомість відстань Фреше (англ. Fréchet Inception Distance, $\mathrm{FID}$) та ядерна відстань (англ. Kernel Inception Distance, $\mathrm{KID}$) оцінюють відстань між розподілами ознак усього набору даних:
                \begin{equation*}
                    \mathrm{FID}(x, y) = \|\mu_x - \mu_y\|^2 + \mathrm{Tr}\left(\Sigma_x + \Sigma_y - 2(\Sigma_x \Sigma_y)^{\frac{1}{2}}\right),
                \end{equation*}
                де $\mu_x, \mu_y$ -- вектори середніх значень ознак; $\Sigma_x, \Sigma_y$ -- матриці коваріації; $\mathrm{Tr}$ -- слід матриці.
                \begin{equation*}
                    \mathrm{KID}(x, y) = {\mathrm{MMD}(x, y)}^2,
                \end{equation*}
                де $\mathrm{MMD}$ -- максимальна середня розбіжність (Maximum Mean Discrepancy) між представленнями ознак.
                Зниження цих показників свідчить про те, що статистика згенерованих зображень наближається до реальних фотографій.
                Оскільки $\mathrm{KID}$ обчислюється ітеративно на випадкових підвибірках,
                результат подається у форматі середнього значення $\mathrm{KID}_{\mathrm{mean}}$ та стандартного відхилення $\mathrm{KID}_{\mathrm{std}}$,
                що характеризує стабільність моделі;

            \item \textbf{обчислювальна та апаратна ефективність:}
                оскільки розроблювана архітектура орієнтована на інтерактивність,
                критичними параметрами є час інференсу, $\mathrm{Time}$, та пікове споживання відеопам'яті (англ. Video Random Access Memory, $\mathrm{VRAM}$).
        \end{itemize}

    % \label{chap:theory}

        % \begin{figure}[h!]
        %     \centering
        %     \includegraphics[width=1\textwidth]{../Preps/StructureScheme.png}
        %     \caption{Структурна схема класифікації методів колоризації}
        %     \label{fig:StructureScheme}
        % \end{figure}

    %     У цьому розділі проводиться більш детальний аналіз сучасних алгоритмів колоризації, що зображені на структурній схемі (рис. \ref{fig:StructureScheme}).
    %     Ми розглянемо еволюцію підходів від класичних оптимізаційних задач до новітніх генеративних моделей.

    %     Основна мета цього розділу -- виявити сильні та слабкі сторони існуючих рішень за такими критеріями: точність відтворення кольорів, обчислювальна складність, можливість керування результатом та часова узгодженість для відеопослідовностей.

    % \section{Алгоритмічні методи колоризації зображень}

    % \section{Методи колоризації на основі нейронних мереж}

    \section{Більш детально про існуючі підходи до колоризації}

        До епохи глибокого навчання колоризація розглядалася переважно як задача оптимізації.
        Класичні алгоритмічні методи можна розділити на дві великі групи: методи на основі втручання користувача та методи перенесення кольору зі зразка.

        \subsection{Методи на основі втручання користувача}
            Цей підхід базується на припущенні, що сусідні пікселі з подібною інтенсивністю повинні мати подібний колір.
            Користувач наносить на чорно-біле зображення кольорові штрихи, які слугують граничними умовами.
            Задача зводиться до мінімізації цільової функції енергії $J$ \cite{2025-survey, 2004-colorization-using-optimization}:

            \begin{equation*}
                J(U) = \sum_{r} \left( U(r) - \sum_{s \in N(r)} w_{rs} U(s) \right)^2,            
            \end{equation*}
            де $U(r)$ -- колір у пікселі $r$, $N(r)$ -- окіл пікселя, а $w_{rs}$ -- ваговий коефіцієнт, що залежить від подібності яскравості пікселів $r$ та $s$.

            Перевагою методу є повний контроль користувача, проте він вимагає значних ручних зусиль і часто дає артефакти на межах об'єктів зі слабким контрастом.

        \subsection{Методи перенесення кольору зі зразка}
            У цьому підході колір переноситься з референсного кольорового зображення на цільове чорно-біле шляхом зіставлення статистик яскравості та текстур.
            
            Ранні роботи \cite{2001-reinhard-color-transfer,2002-welsh-transferring-color} використовували зіставлення гістограм у декорельованому колірному просторі $l\alpha\beta$ (простір Рудермана), що дозволяло маніпулювати колірними каналами незалежно один від одного без виникнення візуальних артефактів.
            
            Хоча цей метод є більш автоматизованим, він критично залежить від вдалого вибору референсного зображення та часто призводить до помилкового забарвлення семантично різних об'єктів, що мають схожу текстуру.

        \subsection{Згорткові нейронні мережі}
            Поява CNN дозволила моделювати складні залежності між формою об'єкта та його ймовірним кольором, використовуючи великі набори даних.

            Перші нейромережеві методи намагалися прямо передбачити значення каналів $ab$ у просторі Lab, мінімізуючи MSE:
            \begin{equation*}
                \mathcal{L}_{L2} = \frac{1}{2} \sum_{h,w} \lVert Y_{gt} - \hat{Y} \rVert_2^2,
            \end{equation*}
            де $Y_{gt}$ -- справжнє зображення, а $\hat{Y}$ -- згенероване.

            Проте, як зазначають Zhang et al. \cite{2016-colorful-image-colorization}, використання MSE призводить до усереднення всіх можливих кольорів, що дає ненасичені результати.

            Для вирішення цієї проблеми задачу було переформульовано як мультимодальну класифікацію, де мережа передбачає розподіл ймовірностей для квантованих відтінків кольору, що значно підвищило яскравість результатів.

        \subsection{Генеративно-змагальні мережі}
            Мережі GAN, зокрема архітектура Pix2Pix \cite{2017-image-to-image-translation}, здійснили прорив у чіткості зображень.
            Генератор $G$ намагається створити реалістичне зображення, щоб обдурити дискримінатор $D$, який вчиться відрізняти справжні кольорові фото від згенерованих.
            Функція втрат cGAN має вигляд:
            \begin{equation*}
                \mathcal{L}_{cGAN}(G, D) = \mathbb{E}_{x, y}[\log D(x, y)] + \mathbb{E}_{x, z}[\log(1 - D(x, G(x, z)))],
            \end{equation*}
            де $x$ -- вхідне чорно-біле зображення, $y$ -- цільове кольорове зображення, а $z$ -- вектор випадкового шуму.
        
            Підходи на кшталт ChromaGAN \cite{2020-chromagan} інтегрують семантичну інформацію у процес генерації, що дозволяє уникнути грубих помилок, наприклад, зеленого неба.

        \subsection{Трансформери та механізми уваги}
            Останні дослідження \cite{2021-colorization-transformer, 2022-unicolor, 2023-ddcolor} показують, що класичні CNN мають локальне рецептивне поле,
            через що вони погано справляються з моделюванням глобальних зв'язків.
            Це часто призводить до семантичних помилок: наприклад, різні частини одного об'єкта,
            перекриті іншими елементами сцени, можуть бути розфарбовані в різні кольори.
            
            Перехід до архітектур на базі трансформерів дозволив вирішити цю проблему.
            Вони використовують механізм глобальної уваги, який дозволяє кожному патчу зображення взаємодіяти з усіма іншими патчами одночасно, враховуючи контекст усієї сцени.
            Це є критично важливим для просторово узгодженого забарвлення великих областей та підтримки єдиної логіки освітлення.
            
            Зокрема, модель Colorization Transformer \cite{2021-colorization-transformer} вперше успішно застосувала цей підхід,
            розбивши зображення на послідовність патчів і розглядаючи задачу колоризації як задачу прогнозування послідовності токенів, подібно до генерації тексту.
            Модель генерує кольори авторегресійно, що забезпечує високу якість, проте робить процес інференсу надзвичайно повільним.

            Значним кроком у напрямку інтерактивності стала модель UniColor \cite{2022-unicolor}.
            Вона запропонувала уніфікований фреймворк, який здатний обробляти мультимодальні умови -- від автоматичної колоризації до генерації на основі текстових описів,
            референсних зображень або точкових підказок. Використання трансформера дозволило ефективно інтегрувати ці різнорідні підказки у єдиний простір репрезентацій.

            Для вирішення проблеми розмиття контурів, яка часто виникає у чистих трансформерів, архітектура DDColor \cite{2023-ddcolor} впровадила систему подвійних декодерів.
            Один декодер відповідає виключно за відновлення просторової роздільної здатності, тобто за піксельний рівень, а інший -- за формування кольорових семантичних ознак, тобто за глобальний контекст, після чого їхні результати зливаються.
            Це дозволило досягти високого рівня фотореалістичності.

            Проте, незважаючи на беззаперечні переваги в якості генерації, використання класичного механізму уваги накладає жорсткі апаратні обмеження.
            Обчислювальна складність глобальної уваги зростає квадратично $\mathcal{O}(N^2)$ зі збільшенням роздільної здатності зображення.
            Це робить такі моделі вкрай вимогливими до обсягу ($\mathrm{VRAM}$) і ускладнює їх використання для інтерактивної колоризації зображень у високій якості в реальному часі.

            Спроби подолати цю проблему за рахунок так званих <<ефективних трансформерів>> призводять до нових компромісів.
            Наприклад, обчислення уваги в локальних вікнах, як-от, в Swin Transformer \cite{2021-swin-transformer},
            або використання математичних лінійних та інших апроксимацій можуть обмежувати здатність моделі до повного глобального моделювання контексту \cite{2021-long-range-arena, 2022-efficient-transformers-survey}.
            У задачі інтерактивної колоризації це стає критичним недоліком, оскільки нейромережа втрачає здатність швидко та безшовно поширювати колір від точкової підказки користувача на віддалені ділянки об'єкта.
            
        \subsection{Дифузійні моделі}
            На відміну від GAN, які генерують результат за один прохід, дифузійні моделі формують зображення ітеративно.
            Математична основа класичних дифузійних моделей базується на двох марковських процесах:
            прямому, який поступово руйнує дані шляхом додавання гаусового шуму,
            та зворотному, де нейронна мережа вчиться крок за кроком відновлювати оригінальний сигнал,
            керуючись вхідним чорно-білим зображенням як просторовою умовою.

            Цей клас моделей здійснив революцію в синтезі зображень завдяки високій стабільності навчання,
            тобто відсутності проблеми колапсу мод, який характерний для GAN, та неперевершеній якості генерації дрібних текстур.
            Знаковою роботою в цій галузі став універсальний фреймворк Palette \cite{2022-palette-image-to-image-diffusion-models},
            який продемонстрував, що умовна дифузійна модель здатна перевершити спеціалізовані GAN-архітектури у задачі колоризації без додаткового налаштування під специфіку предметної області. 

            Поряд із класичним підходом активно розвиваються альтернативні формулювання генеративного процесу.
            Наприклад, метод Cold Diffusion \cite{2023-cold-diffusion} довів можливість інверсії довільних детермінованих перетворень зображень без використання чистого гаусового шуму,
            що розширює теоретичні межі застосування дифузійних алгоритмів.

            Проте фундаментальним недоліком дифузійних моделей залишається їхня надзвичайно висока обчислювальна вартість на етапі інференсу.
            Оскільки для створення фінального зображення модель повинна виконати десятки або сотні послідовних ітерацій денойзингу, час генерації зростає на порядки порівняно з моделями прямого проходу.
            У контексті керованої інтерактивної колоризації, де фундаментальною вимогою є миттєвий візуальний зворотний зв'язок на кожну нову підказку користувача, використання чистих дифузійних моделей є недоцільним через руйнування користувацького досвіду.
        
        \subsection{Селективні моделі простору станів та архітектура Mamba}
            Новітнім класом нейромережевих архітектур, що дозволяє подолати квадратичну складність трансформерів та повільність дифузійних моделей,
            є SSM, зокрема архітектура Mamba \cite{2023-mamba}.
            Вони здатні моделювати наскрізні глобальні контекстні зв'язки з суворою лінійною обчислювальною складністю $\mathcal{O}(L)$,
            що робить їх ідеальними для систем реального часу.

            Математично базова система SSM відображає вхідний одновимірний сигнал $x(t) \in \mathbb{R}$ у прихований стан $h(t) \in \mathbb{R}^N$ через лінійні диференціальні рівняння:
            \begin{align*}
                h'(t) &= A \cdot h(t) + B \cdot x(t), \\
                y(t) &= C \cdot h(t) + D \cdot x(t),
            \end{align*}
            де $A, B, C, D$ -- матриці параметрів системи.
            При цьому матриця $A$ характеризує внутрішню динаміку системи та керує еволюцією прихованого стану, тобто відповідає за довготривалу пам'ять,
            матриця $D$ описує прямий зв'язок між входом і виходом, виконуючи роль класичного пропускного з'єднання,
            а матриці $B$ та $C$ визначають коефіцієнти взаємодії прихованого стану із вхідним та вихідним сигналами відповідно.

            Для обробки дискретних послідовностей безперервні параметри трансформують у дискретні ($\overline{A}, \overline{B}$) із використанням кроку дискретизації $\Delta$. Зазвичай для цього застосовують метод фіксації нульового порядку (англ. Zero-Order Hold):
            \begin{align*}
                \overline{A} &= \exp(\Delta A), \\
                \overline{B} &= (\Delta A)^{-1}(\exp(\Delta A) - I) \cdot \Delta B,
            \end{align*}
            що дозволяє переписати систему у дискретному вигляді: $h_k = \overline{A} \cdot h_{k-1} + \overline{B} \cdot x_k$, $y_k = C \cdot h_k + D \cdot x_k$.

            Головна інновація архітектури Mamba полягає у селективності її внутрішніх механізмів:
            параметри $B$, $C$ та крок дискретизації $\Delta$ стають динамічними функціями від вхідного сигналу,
            на відміну від статичних матриць $A$ та $D$.
            Завдяки такій дискретизації параметрів та використанню апаратно-оптимізованого алгоритму паралельного сканування в пам'яті графічного процесора,
            модель динамічно фільтрує нерелевантну інформацію та ефективно обробляє довгі послідовності без обчислювальних обмежень механізму уваги.

            Однак пряме застосування оригінальної одновимірної Mamba до двовимірних зображень призводить до проблеми втрати локального контексту,
            оскільки розгортання матриці пікселів у вектор розриває їхню просторову суміжність по вертикалі.
            Для вирішення цієї проблеми візуальні адаптації на кшталт Vision Mamba \cite{2024-vim} та новітні моделі спектральної трансляції \cite{2024-colormamba} використовують механізми двонаправленого або перехресного сканування.
            Це дозволяє моделі агрегувати ознаки з різних напрямків, повністю відновлюючи просторове розуміння сцени.

            У контексті керованої колоризації використання Mamba відкриває унікальну можливість:
            здатність моделі миттєво та узгоджено поширювати колірну інформацію від точкової підказки користувача на весь об'єкт завдяки глобальному рецептивному полю,
            уникаючи при цьому експоненційного зростання споживання відеопам'яті.
            Саме цей синергетичний баланс обчислювальної ефективності та просторової експресивності робить SSM найбільш перспективним фундаментом для розробки сучасних інтерактивних систем розфарбовування.

    \section{Порівняння різних методів колоризації}

        Комплексне порівняння розглянутих методів (табл. \ref{tab:comparison}) базується на чотирьох ключових критеріях (Рис. \ref{fig:comparison_criteria}).
        
        \begin{figure}[ht]
            \centering
            \begin{tikzpicture}[
                node distance=0.5cm and 0.5cm,
                every node/.style={
                    font=\small,
                    align=center
                },
                main/.style={
                    draw,
                    rounded corners,
                    thick,
                    fill=gray!15,
                    minimum width=4.5cm,
                    minimum height=1.2cm
                },
                block/.style={
                    draw,
                    rounded corners,
                    thick,
                    fill=blue!10,
                    text width=4.2cm,
                    minimum height=2.2cm
                }
            ]

            % Центральний блок
            \node[main] (center) {Ключові критерії\\ аналізу методів};

            % Верхній лівий
            \node[block, above left=of center] (quality) {
                \textbf{Якість}\\
                Фотореалістичність,\\
                природність та\\
                насиченість кольорів
            };

            % Верхній правий
            \node[block, above right=of center] (consistency) {
                \textbf{Узгодженість}\\
                Просторова зв'язність\\
                кольору та відсутність\\
                артефактів
            };

            % Нижній лівий
            \node[block, below left=of center] (control) {
                \textbf{Керованість}\\
                Точність інтерпретації\\
                та інтеграції\\
                підказок користувача
            };

            % Нижній правий
            \node[block, below right=of center] (speed) {
                \textbf{Швидкість}\\
                Обчислювальна\\
                складність та\\
                час інференсу
            };

            % Лінії
            \draw[thick] (center) -- (quality);
            \draw[thick] (center) -- (consistency);
            \draw[thick] (center) -- (control);
            \draw[thick] (center) -- (speed);

            \end{tikzpicture}

            \caption{Ключові критерії порівняння методів колоризації}
            \label{fig:comparison_criteria}
        \end{figure}

        \begin{table}[ht!]
            \centering
            \caption{Порівняльний аналіз методів колоризації зображень}
            \label{tab:comparison}
            \begin{tblr}{
                width = \textwidth,
                colspec = {X[2,l,t] X[2,l,t] X[3,l,t] X[3,l,t]},
                row{1} = {font=\bfseries, halign=c},
                hlines,
                vlines,
            }
                Метод / Архітектура    & Представники                                                                                                                  & Переваги                                                                                                    & Недоліки                                                                                                                 \\
                Класичні оптимізаційні & Levin et al. \cite{2004-colorization-using-optimization}, Welsh \cite{2002-welsh-transferring-color}                          & Точний локальний контроль користувача, відсутність семантичних <<галюцинацій>>                              & Обчислювально повільні, вимагають значної ручної розмітки, пасують перед складними текстурами                            \\
                CNN                    & Zhang et al. \cite{2016-colorful-image-colorization}, Larsson \cite{2016-learning-representations-for-automatic-colorization} & Висока швидкість інференсу, стабільність процесу навчання                                                   & Обмежене рецептивне поле; схильність до генерації усереднених, тьмяних відтінків при мінімізації $\mathrm{MSE}$          \\
                GAN                    & Pix2Pix \cite{2017-image-to-image-translation}, ChromaGAN \cite{2020-chromagan}                                               & Висока чіткість контурів, насичені та візуально привабливі кольори                                          & Нестабільність змагального навчання, схильність до появи колірних артефактів та плям                                     \\
                Transformers           & ColTran \cite{2021-colorization-transformer}, UniColor \cite{2022-unicolor}, DDColor \cite{2023-ddcolor}                      & Моделювання глобального контексту всієї сцени, гнучка мультимодальна керованість                            & Квадратична обчислювальна складність $\mathcal{O}(N^2)$, критично високі вимоги до обсягу відеопам'яті ($\mathrm{VRAM}$) \\
                Diffusion Models       & Palette \cite{2022-palette-image-to-image-diffusion-models}, Cold Diffusion \cite{2023-cold-diffusion}                        & Найвища якість генерації (SOTA), бездоганний фотореалізм текстур                                            & Екстремально низька швидкість інференсу через сотні послідовних ітерацій денойзингу                                      \\
                Mamba SSM              & Vision Mamba \cite{2024-vim}, ColorMamba \cite{2024-colormamba}                                                               & Суворо лінійна складність $\mathcal{O}(N)$, стабільне споживання $\mathrm{VRAM}$, глобальне рецептивне поле & Чутливість до одновимірного розгортання 2D-даних, ризик втрати або розмиття локальних контурів об'єктів                  \\
            \end{tblr}
        \end{table}

        Проведений аналіз демонструє наявність компромісу між обчислювальною лінійністю та збереженням локальних просторових контурів у сучасних послідовних моделях.
        Чисті трансформери забезпечують високу якість зв'язності, але стають апаратно непридатними для обробки великих зображень у реальному часі, тоді як SSM пропонують бажану лінійну ефективність $\mathcal{O}(N)$,
        але потребують додаткових механізмів локального посилення контурів. 

        Таким чином, гібридні підходи, що поєднують генеративну силу та швидкість селективних лінійних моделей простору станів із точністю локального вилучення ознак за допомогою ієрархічних згорткових декодерів,
        виглядають найбільш обґрунтованими та перспективними для вирішення задачі керованої інтерактивної колоризації.
        Саме розробці та дослідженню такої гібридної архітектури присвячено наступні розділи цієї роботи.

\chapconclude{\ref{chap:review_theory}}

    У цьому розділі було проведено огляд теоретичних основ та сучасних методів колоризації зображень.
    Було показано, що еволюція методів пройшла шлях від простих евристик до складних генеративних моделей, таких як GAN, Transformers, дифузійні фреймворки та селективні моделі простору станів.

    Попри значні успіхи безумовної колоризації у досягненні реалістичності результатів, залишаються відкритими питання обчислювальної складності моделей під час роботи із зображеннями високої роздільної здатності, керованості процесом колоризації та забезпечення багатоваріантності результатів.
    Зокрема, гібридні підходи, мультимодальні підходи та уніфіковані дифузійні й рекурентні структури є найбільш актуальним напрямком досліджень.

    Порівняння нейромережевих підходів показало, що хоча архітектури GAN забезпечують швидкий інференс, вони часто страждають від нестабільності навчання та появи колірних артефактів.
    Дифузійні моделі дозволяють досягти найвищої якості та фотореалізму текстур, проте їхнє практичне застосування жорстко обмежене тривалим часом інференсу через ітеративний характер процесу генерації.
    У свою чергу, підходи на базі трансформерів ефективно моделюють глобальний контекст усієї сцени, однак квадратичне зростання обчислювальних витрат $\mathcal{O}(N^2)$ суттєво ускладнює їх масштабування для зображень високої роздільної здатності.
    
    У наступному розділі буде запропоновано власний підхід до вирішення задачі колоризації зображень, що базується на гібридизації архітектури U-Net та Mamba SSM з використанням спільних ваг.
    Це дозволить ефективно поєднати вилучення низькорівневих просторових ознак, для боротьби з color bleeding, із лінійним моделюванням глобального контексту та суттєвим збереженням пам'яті графічного процесора.

    % Наприкінці кожного розділу ви повинні навести коротенькі підсумки по його 
    % результатах. Зокрема, для оглядового розділу в якості висновків необхідно 
    % зазначити, які задачі у даній тематиці вже були розв'язані, а саме 
    % поставлена вами задача розв'язана не була (або розв'язана погано), тому у 
    % наступних розділах ви її й розв'язуєте.

    % До другого розділу також краще написати малесенький вступ. Зокрема, це 
    % збільшує загальний об'єм роботи та покращує її читабельність.


    % У другому розділі необхідно наводити розв'язання поставленої перед вами 
    % задачі у теоретичному або аналітичному сенсі (хоча, звісно, все залежить 
    % від того, яка саме задача перед вами поставлена).

    % Бажано, щоб кожен пункт завдань, окреслених у вступі, відповідав певному 
    % розділу або підрозділу у дипломній роботі.

    % Посилання до роботи рекомендовано робити за допомогою підсистеми BibTex. 
    % Це максимально зручно, повірте мені. Усі притомні світові журнали надають 
    % посилання на свої статті у вигляді bib-файлів, з яких просто треба 
    % перенести відомості у свій файл. Існують інструменти автоматизованого створення bib-файлів, наприклад, JabRef.
    % Якщо у вас нема автоматизованого механізму, то посилання на джерела у 
    % bib-файлах легко створюються вручну. Головне -- не забувати чотири простих 
    % правила:

    % 1) Імена авторів необхідно подавати у вигляді <<Прізвище, Ім'я>> та 
    % розділяти ключовим словом <<and>>, наприклад:

    % \texttt{author = ''Яковлєв, Сергій Володимирович and Дамблдор, Альбус Персиваль Вульфрік Брайан''}

    % 2) Після усіх полів запису необхідно ставити кому, окрім останнього.

    % 3) Джерела у переліку посилань з'являться тільки після того, як ви зробите 
    % посилання на нього у тексті, наприклад, див. роботи~\cite{ct1, ct2}.

    % Наприкінці розділу знову наводяться коротенькі підсумки.